
\documentclass{article}
\usepackage{colortbl}
\usepackage{makecell}
\usepackage{multirow}
\usepackage{supertabular}

\begin{document}

\newcounter{utterance}

\centering \large Interaction Transcript for game `dond', experiment `coop\_it', episode 19 with claude{-}sonnet{-}4{-}20250514{-}no{-}reasoning{-}t1.0.
\vspace{24pt}

{ \footnotesize  \setcounter{utterance}{1}
\setlength{\tabcolsep}{0pt}
\begin{supertabular}{c@{$\;$}|p{.15\linewidth}@{}p{.15\linewidth}p{.15\linewidth}p{.15\linewidth}p{.15\linewidth}p{.15\linewidth}}
    \# & $\;$A & \multicolumn{4}{c}{Game Master} & $\;\:$B\\
    \hline

    \theutterance \stepcounter{utterance}  
    & & \multicolumn{4}{p{0.6\linewidth}}{
        \cellcolor[rgb]{0.9,0.9,0.9}{
            \makecell[{{p{\linewidth}}}]{
                \texttt{\tiny{[P1$\langle$GM]}}
                \texttt{State giocando a un gioco di negoziazione cooperativa in cui dovete accordarvi su come dividere una serie di oggetti tra voi e un altro giocatore.} \\
\\ 
\texttt{Regole:} \\
\texttt{(a) A Lei e all'altro giocatore viene dato un insieme di oggetti. Ognuno di voi riceve anche una funzione di valore segreta, che rappresenta il valore di ciascun tipo di oggetto.} \\
\texttt{(b) Si scambiano messaggi con l'altro giocatore per concordare chi si aggiudica gli oggetti. Potete inviare un massimo di 5 messaggi ciascuno, oppure terminare in anticipo facendo una proposta segreta in qualsiasi momento.} \\
\texttt{(c) A ciascuno di voi viene chiesto di inviare una proposta segreta indicando gli oggetti che desiderate, formattata tra parentesi quadre come segue: "{[}Proposta: <numero> <nome oggetto>, <numero> <nome oggetto>, <...>{]}".} \\
\texttt{(d) Se le vostre proposte sono complementari, cioè ci sono abbastanza oggetti per soddisfare entrambe le proposte, a ciascun giocatore viene assegnato un punteggio basato sulla somma dei valori degli oggetti ricevuti. In caso contrario, entrambi ricevono zero punti.} \\
\texttt{(e) L'obiettivo è massimizzare la somma del proprio punteggio e di quello dell'altro giocatore.} \\
\\ 
\texttt{Cominciamo.} \\
\\ 
\texttt{L'insieme degli oggetti disponibili è:} \\
\\ 
\texttt{2 ganci, 1 fiore, 1 scheda, 1 gomma, 1 spugna} \\
\\ 
\texttt{La funzione valore segreta è:} \\
\\ 
\texttt{gancio: 3, fiore: 3, scheda: 0, gomma: 1, spugna: 0} \\
\\ 
\texttt{IMPORTANTE: i vostri messaggi, a meno che non si tratti di una proposta segreta, vengono trasmessi direttamente all'altro giocatore, quindi non includete alcuna risposta alle regole o testo di annuncio del vostro messaggio. Per fare una proposta segreta, utilizzate il formato indicato. Non utilizzare delle parentesi quadre quando si comunica all'altro giocatore, altrimenti verrà interpretata come una proposta segreta.} \\
            }
        }
    }
    & \\ \\

    \theutterance \stepcounter{utterance}  
    & \multicolumn{4}{p{0.6\linewidth}}{
        \cellcolor[rgb]{0.9,0.9,0.9}{
            \makecell[{{p{\linewidth}}}]{
                \texttt{\tiny{[P1$\rangle$GM]}}
                \texttt{Ciao! Vedo che abbiamo 2 ganci, 1 fiore, 1 scheda, 1 gomma e 1 spugna da dividere tra noi. Per trovare la divisione migliore possibile, dovremmo capire quali oggetti sono più preziosi per ciascuno di noi. } \\
\\ 
\texttt{Potresti dirmi quali oggetti consideri più importanti o di maggior valore? Io ho delle preferenze specifiche, ma vorrei sentire le tue per trovare un accordo che funzioni bene per entrambi.} \\
            }
        }
    }
    & & \\ \\

    \theutterance \stepcounter{utterance}  
    & & & \multicolumn{2}{p{0.3\linewidth}}{
        \cellcolor[rgb]{0.9,0.9,0.9}{
            \makecell[{{p{\linewidth}}}]{
                \texttt{\tiny{[GM$|$GM]}}
                \texttt{continue} \\
            }
        }
    }
    & & \\ \\

    \theutterance \stepcounter{utterance}  
    & & & \multicolumn{2}{p{0.3\linewidth}}{
        \cellcolor[rgb]{0.9,0.9,0.9}{
            \makecell[{{p{\linewidth}}}]{
                \texttt{\tiny{[GM$|$GM]}}
                \texttt{Ciao! Vedo che abbiamo 2 ganci, 1 fiore, 1 scheda, 1 gomma e 1 spugna da dividere tra noi. Per trovare la divisione migliore possibile, dovremmo capire quali oggetti sono più preziosi per ciascuno di noi. } \\
\\ 
\texttt{Potresti dirmi quali oggetti consideri più importanti o di maggior valore? Io ho delle preferenze specifiche, ma vorrei sentire le tue per trovare un accordo che funzioni bene per entrambi.} \\
            }
        }
    }
    & & \\ \\

    \theutterance \stepcounter{utterance}  
    & & \multicolumn{4}{p{0.6\linewidth}}{
        \cellcolor[rgb]{0.9,0.9,0.9}{
            \makecell[{{p{\linewidth}}}]{
                \texttt{\tiny{[GM$\rangle$P2]}}
                \texttt{State giocando a un gioco di negoziazione cooperativa in cui dovete accordarvi su come dividere una serie di oggetti tra voi e un altro giocatore.} \\
\\ 
\texttt{Regole:} \\
\texttt{(a) A Lei e all'altro giocatore viene dato un insieme di oggetti. Ognuno di voi riceve anche una funzione di valore segreta, che rappresenta il valore di ciascun tipo di oggetto.} \\
\texttt{(b) Si scambiano messaggi con l'altro giocatore per concordare chi si aggiudica gli oggetti. Potete inviare un massimo di 5 messaggi ciascuno, oppure terminare in anticipo facendo una proposta segreta in qualsiasi momento.} \\
\texttt{(c) A ciascuno di voi viene chiesto di inviare una proposta segreta indicando gli oggetti che desiderate, formattata tra parentesi quadre come segue: "{[}Proposta: <numero> <nome oggetto>, <numero> <nome oggetto>, <...>{]}".} \\
\texttt{(d) Se le vostre proposte sono complementari, cioè ci sono abbastanza oggetti per soddisfare entrambe le proposte, a ciascun giocatore viene assegnato un punteggio basato sulla somma dei valori degli oggetti ricevuti. In caso contrario, entrambi ricevono zero punti.} \\
\texttt{(e) L'obiettivo è massimizzare la somma del proprio punteggio e di quello dell'altro giocatore.} \\
\\ 
\texttt{Cominciamo.} \\
\\ 
\texttt{L'insieme degli oggetti disponibili è:} \\
\\ 
\texttt{2 ganci, 1 fiore, 1 scheda, 1 gomma, 1 spugna} \\
\\ 
\texttt{La funzione valore segreta è:} \\
\\ 
\texttt{gancio: 0, fiore: 2, scheda: 1, gomma: 4, spugna: 3} \\
\\ 
\texttt{IMPORTANTE: i vostri messaggi, a meno che non si tratti di una proposta segreta, vengono trasmessi direttamente all'altro giocatore, quindi non includete alcuna risposta alle regole o testo di annuncio del vostro messaggio. Per fare una proposta segreta, utilizzate il formato indicato. Non utilizzare delle parentesi quadre quando si comunica all'altro giocatore, altrimenti verrà interpretata come una proposta segreta.} \\
\\ 
\\ 
\texttt{Ciao! Vedo che abbiamo 2 ganci, 1 fiore, 1 scheda, 1 gomma e 1 spugna da dividere tra noi. Per trovare la divisione migliore possibile, dovremmo capire quali oggetti sono più preziosi per ciascuno di noi. } \\
\\ 
\texttt{Potresti dirmi quali oggetti consideri più importanti o di maggior valore? Io ho delle preferenze specifiche, ma vorrei sentire le tue per trovare un accordo che funzioni bene per entrambi.} \\
            }
        }
    }
    & \\ \\

    \theutterance \stepcounter{utterance}  
    & & & \multicolumn{4}{p{0.6\linewidth}}{
        \cellcolor[rgb]{0.9,0.9,0.9}{
            \makecell[{{p{\linewidth}}}]{
                \texttt{\tiny{[GM$\langle$P2]}}
                \texttt{Ciao! Sì, è una buona strategia condividere le nostre preferenze per trovare la divisione ottimale.} \\
\\ 
\texttt{Dal mio punto di vista, gli oggetti in ordine di importanza sono: gomma (molto importante), spugna (abbastanza importante), fiore (moderatamente importante), scheda (poco importante), ganci (non importanti per me).} \\
\\ 
\texttt{E tu? Quali oggetti consideri più preziosi? Se riusciamo a capire dove le nostre preferenze si complementano, dovremmo riuscire a trovare una divisione che massimizzi il valore totale per entrambi.} \\
            }
        }
    }
     \\ \\

    \theutterance \stepcounter{utterance}  
    & & & \multicolumn{2}{p{0.3\linewidth}}{
        \cellcolor[rgb]{0.9,0.9,0.9}{
            \makecell[{{p{\linewidth}}}]{
                \texttt{\tiny{[GM$|$GM]}}
                \texttt{continue} \\
            }
        }
    }
    & & \\ \\

    \theutterance \stepcounter{utterance}  
    & & & \multicolumn{2}{p{0.3\linewidth}}{
        \cellcolor[rgb]{0.9,0.9,0.9}{
            \makecell[{{p{\linewidth}}}]{
                \texttt{\tiny{[GM$|$GM]}}
                \texttt{Ciao! Sì, è una buona strategia condividere le nostre preferenze per trovare la divisione ottimale.} \\
\\ 
\texttt{Dal mio punto di vista, gli oggetti in ordine di importanza sono: gomma (molto importante), spugna (abbastanza importante), fiore (moderatamente importante), scheda (poco importante), ganci (non importanti per me).} \\
\\ 
\texttt{E tu? Quali oggetti consideri più preziosi? Se riusciamo a capire dove le nostre preferenze si complementano, dovremmo riuscire a trovare una divisione che massimizzi il valore totale per entrambi.} \\
            }
        }
    }
    & & \\ \\

    \theutterance \stepcounter{utterance}  
    & & \multicolumn{4}{p{0.6\linewidth}}{
        \cellcolor[rgb]{0.9,0.9,0.9}{
            \makecell[{{p{\linewidth}}}]{
                \texttt{\tiny{[P1$\langle$GM]}}
                \texttt{Ciao! Sì, è una buona strategia condividere le nostre preferenze per trovare la divisione ottimale.} \\
\\ 
\texttt{Dal mio punto di vista, gli oggetti in ordine di importanza sono: gomma (molto importante), spugna (abbastanza importante), fiore (moderatamente importante), scheda (poco importante), ganci (non importanti per me).} \\
\\ 
\texttt{E tu? Quali oggetti consideri più preziosi? Se riusciamo a capire dove le nostre preferenze si complementano, dovremmo riuscire a trovare una divisione che massimizzi il valore totale per entrambi.} \\
            }
        }
    }
    & \\ \\

    \theutterance \stepcounter{utterance}  
    & \multicolumn{4}{p{0.6\linewidth}}{
        \cellcolor[rgb]{0.9,0.9,0.9}{
            \makecell[{{p{\linewidth}}}]{
                \texttt{\tiny{[P1$\rangle$GM]}}
                \texttt{Perfetto! Le nostre preferenze si complementano molto bene. Per me l'ordine di importanza è: ganci e fiore (molto importanti), gomma (poco importante), mentre scheda e spugna non hanno valore per me.} \\
\\ 
\texttt{Quindi sembra che possiamo creare una divisione win{-}win:} \\
\texttt{{-} Tu prendi: gomma, spugna, scheda (i tuoi oggetti più importanti)} \\
\texttt{{-} Io prendo: 2 ganci, fiore (i miei oggetti più importanti)} \\
\\ 
\texttt{Questa divisione dovrebbe massimizzare il valore totale perché ognuno di noi ottiene gli oggetti che considera più preziosi. Che ne pensi di questa proposta?} \\
            }
        }
    }
    & & \\ \\

    \theutterance \stepcounter{utterance}  
    & & & \multicolumn{2}{p{0.3\linewidth}}{
        \cellcolor[rgb]{0.9,0.9,0.9}{
            \makecell[{{p{\linewidth}}}]{
                \texttt{\tiny{[GM$|$GM]}}
                \texttt{continue} \\
            }
        }
    }
    & & \\ \\

    \theutterance \stepcounter{utterance}  
    & & & \multicolumn{2}{p{0.3\linewidth}}{
        \cellcolor[rgb]{0.9,0.9,0.9}{
            \makecell[{{p{\linewidth}}}]{
                \texttt{\tiny{[GM$|$GM]}}
                \texttt{Perfetto! Le nostre preferenze si complementano molto bene. Per me l'ordine di importanza è: ganci e fiore (molto importanti), gomma (poco importante), mentre scheda e spugna non hanno valore per me.} \\
\\ 
\texttt{Quindi sembra che possiamo creare una divisione win{-}win:} \\
\texttt{{-} Tu prendi: gomma, spugna, scheda (i tuoi oggetti più importanti)} \\
\texttt{{-} Io prendo: 2 ganci, fiore (i miei oggetti più importanti)} \\
\\ 
\texttt{Questa divisione dovrebbe massimizzare il valore totale perché ognuno di noi ottiene gli oggetti che considera più preziosi. Che ne pensi di questa proposta?} \\
            }
        }
    }
    & & \\ \\

    \theutterance \stepcounter{utterance}  
    & & \multicolumn{4}{p{0.6\linewidth}}{
        \cellcolor[rgb]{0.9,0.9,0.9}{
            \makecell[{{p{\linewidth}}}]{
                \texttt{\tiny{[GM$\rangle$P2]}}
                \texttt{Perfetto! Le nostre preferenze si complementano molto bene. Per me l'ordine di importanza è: ganci e fiore (molto importanti), gomma (poco importante), mentre scheda e spugna non hanno valore per me.} \\
\\ 
\texttt{Quindi sembra che possiamo creare una divisione win{-}win:} \\
\texttt{{-} Tu prendi: gomma, spugna, scheda (i tuoi oggetti più importanti)} \\
\texttt{{-} Io prendo: 2 ganci, fiore (i miei oggetti più importanti)} \\
\\ 
\texttt{Questa divisione dovrebbe massimizzare il valore totale perché ognuno di noi ottiene gli oggetti che considera più preziosi. Che ne pensi di questa proposta?} \\
            }
        }
    }
    & \\ \\

    \theutterance \stepcounter{utterance}  
    & & & \multicolumn{4}{p{0.6\linewidth}}{
        \cellcolor[rgb]{0.9,0.9,0.9}{
            \makecell[{{p{\linewidth}}}]{
                \texttt{\tiny{[GM$\langle$P2]}}
                \texttt{Perfetto! Questa divisione mi sembra ottimale. Effettivamente ognuno di noi ottiene gli oggetti che considera più preziosi:} \\
\\ 
\texttt{{-} Tu ottieni i 2 ganci e il fiore, che sono i tuoi oggetti più importanti} \\
\texttt{{-} Io ottengo gomma, spugna e scheda, che sono esattamente gli oggetti che valuto di più} \\
\\ 
\texttt{È una situazione ideale dove non c'è conflitto sugli oggetti desiderati. Procediamo con questa divisione!} \\
\\ 
\texttt{{[}Proposta: 1 gomma, 1 spugna, 1 scheda{]}} \\
            }
        }
    }
     \\ \\

    \theutterance \stepcounter{utterance}  
    & & & \multicolumn{2}{p{0.3\linewidth}}{
        \cellcolor[rgb]{0.9,0.9,0.9}{
            \makecell[{{p{\linewidth}}}]{
                \texttt{\tiny{[GM$|$GM]}}
                \texttt{proposal} \\
            }
        }
    }
    & & \\ \\

    \theutterance \stepcounter{utterance}  
    & & & \multicolumn{2}{p{0.3\linewidth}}{
        \cellcolor[rgb]{0.9,0.9,0.9}{
            \makecell[{{p{\linewidth}}}]{
                \texttt{\tiny{[GM$|$GM]}}
                [0, 0, 1, 1, 1]
            }
        }
    }
    & & \\ \\

    \theutterance \stepcounter{utterance}  
    & & \multicolumn{4}{p{0.6\linewidth}}{
        \cellcolor[rgb]{0.9,0.9,0.9}{
            \makecell[{{p{\linewidth}}}]{
                \texttt{\tiny{[P1$\langle$GM]}}
                \texttt{L'altro giocatore ha presentato una proposta segreta. Ora tocca a Lei presentare una sua proposta.} \\
            }
        }
    }
    & \\ \\

    \theutterance \stepcounter{utterance}  
    & \multicolumn{4}{p{0.6\linewidth}}{
        \cellcolor[rgb]{0.9,0.9,0.9}{
            \makecell[{{p{\linewidth}}}]{
                \texttt{\tiny{[P1$\rangle$GM]}}
                \texttt{{[}Proposta: 2 ganci, 1 fiore{]}} \\
            }
        }
    }
    & & \\ \\

    \theutterance \stepcounter{utterance}  
    & & & \multicolumn{2}{p{0.3\linewidth}}{
        \cellcolor[rgb]{0.9,0.9,0.9}{
            \makecell[{{p{\linewidth}}}]{
                \texttt{\tiny{[GM$|$GM]}}
                \texttt{proposal} \\
            }
        }
    }
    & & \\ \\

    \theutterance \stepcounter{utterance}  
    & & & \multicolumn{2}{p{0.3\linewidth}}{
        \cellcolor[rgb]{0.9,0.9,0.9}{
            \makecell[{{p{\linewidth}}}]{
                \texttt{\tiny{[GM$|$GM]}}
                [2, 1, 0, 0, 0]
            }
        }
    }
    & & \\ \\

    \theutterance \stepcounter{utterance}  
    & & & \multicolumn{2}{p{0.3\linewidth}}{
        \cellcolor[rgb]{0.9,0.9,0.9}{
            \makecell[{{p{\linewidth}}}]{
                \texttt{\tiny{[GM$|$GM]}}
                [[2, 1, 0, 0, 0], [0, 0, 1, 1, 1]]
            }
        }
    }
    & & \\ \\

\end{supertabular}
}

\end{document}
